% A representative NeurIPS-style paper. Written for ByeTex test purposes;
% uses the documented neurips_2024.sty patterns. To compile with pdflatex,
% obtain neurips_2024.sty from https://neurips.cc/Conferences/2024/PaperInformation/StyleFiles
\documentclass{article}

\PassOptionsToPackage{numbers, compress}{natbib}
\usepackage[final]{neurips_2024}

\usepackage[utf8]{inputenc}
\usepackage[T1]{fontenc}
\usepackage{hyperref}
\usepackage{url}
\usepackage{booktabs}
\usepackage{amsfonts}
\usepackage{nicefrac}
\usepackage{microtype}
\usepackage{graphicx}

\title{A Generic NeurIPS Paper Template}

\author{%
  Alice Anderson\thanks{Equal contribution. Correspondence to: alice@example.edu.} \\
  Department of Computer Science \\
  Example University \\
  Springfield, USA \\
  \texttt{alice@example.edu} \\
  \And
  Bob Brown \\
  Department of Statistics \\
  Example University \\
  Springfield, USA \\
  \texttt{bob@example.edu}
}

\begin{document}

\maketitle

\begin{abstract}
The abstract paragraph should be indented \nicefrac{1}{2}~inch (3~picas) on
both the left- and right-hand margins. Use 10~point type, with a vertical
spacing (leading) of 11~points. The word \textbf{Abstract} must be centered,
bold, and in point size 12. Two line spaces precede the abstract. The abstract
must be limited to one paragraph.
\end{abstract}

\section{Introduction}

A neural network with parameters $\theta \in \mathbb{R}^d$ minimizes a loss
function $\mathcal{L}(\theta)$ via stochastic gradient descent:
\begin{equation}
\theta_{t+1} = \theta_t - \eta_t \nabla \mathcal{L}(\theta_t).
\label{eq:sgd}
\end{equation}

\section{Background}

\subsection{Notation}

We use lowercase letters for scalars ($x \in \mathbb{R}$), bold lowercase for
vectors ($\mathbf{x} \in \mathbb{R}^d$), and bold uppercase for matrices
($\mathbf{A} \in \mathbb{R}^{m \times n}$). The Frobenius norm is
$\|\mathbf{A}\|_F = \sqrt{\sum_{i,j} A_{ij}^2}$.

\subsection{Related work}

Prior approaches~\citep{einstein,dirac} formulated this problem differently.
\citet{bohr} showed that the canonical form is optimal under mild assumptions.

\section{Method}

We propose a two-stage procedure. First, we align by minimizing
\begin{align}
\mathcal{L}_{\text{align}}(\theta) &= \frac{1}{N} \sum_{i=1}^{N} \| f_\theta(x_i) - y_i \|_2^2 \\
&\quad + \lambda \| \theta \|_2^2.
\end{align}
Second, we fine-tune with the standard cross-entropy loss.

\subsection{Implementation details}

We trained for $200$ epochs with learning rate $\eta = 10^{-3}$. See
Table~\ref{tab:results} for the final numbers.

\begin{table}[h]
  \caption{Final results across three random seeds.}
  \label{tab:results}
  \centering
  \begin{tabular}{lrr}
    \toprule
    Method     & Accuracy & Wall time (h) \\
    \midrule
    Baseline   & 81.2     & 12.3 \\
    Ours       & 93.4     & 14.7 \\
    \bottomrule
  \end{tabular}
\end{table}

\section{Discussion}

Equation~\eqref{eq:sgd} is the canonical SGD update. Our refinement applies
when the gradient is sparse.

\section*{Acknowledgments}

We thank the anonymous reviewers and our funding agency for support.

\bibliographystyle{plainnat}
\bibliography{refs}

\end{document}
